% Folder 78, batch 04 (archivists' pages 61–80).
% Pass: Opus 5.5 (claude-opus-5-5), 2026-10-02 — first pass, unchecked against the pages by a human.
%
\documentclass[11pt,a4paper]{article}
\input{../preamble/grothendieck}

\folder{78}
\batch{4}
\pages{61}{80}
\foldertitle{[Polygones réguliers et polynômes cyclotomiques] : notes manuscrites (s.d.).}
\dating{[à partir de 1977]}
\watermark{Édition de démonstration}

\begin{document}

\page{61}
\note{La page s'ouvre sur un point c) : les points a) et b) de cette énumération sont sur les pages précédentes, hors de ce lot.}
c) Pour tout $\omega \in \mathrm{Or}(\Pi)$, $\exists!\, u_\omega \in G^{+}$ tel que $u_\omega | R^{+}_{\omega} = u | R^{+}_{\omega}$, et on a \struck{\ill{}}
\[
u_{-\omega} = u_\omega^{-1}
\]
\struck{Donc} De plus, $u_\omega$ est un générateur de $G^{+}$, et si $n \geq 3$, $\omega \mapsto u_\omega$ est une injection
\[
\mathrm{Or}(\Pi) \hookrightarrow \mathrm{Génér.}(G^{+})
\]
\marginal{flèches tracées dans la marge gauche, pointant vers le point d)}

d) Le foncteur
\[
\Pi \longmapsto \bigl(\mathbf{Z} G^{+},\ \{u_\omega, u_{-\omega}\} \in G^{+},\ S\bigr)
\]
\note{le premier terme se lit « $\mathbf{Z}$ » accolé à $G^{+}$ (peut-être $\mathbf{Z} \simeq G^{+}$ pour $n=\infty$) ; au-dessus de la paire, renvoi « $\Omega$ » ; sous la formule, ajout « ($3 \leq n \leq +\infty$) » rattaché à « $n$-polygones ».}
des \add{$n$-}polygones vers les triples \uncertain{formés} d'un groupe cyclique d'ordre $n$ $G^{+}$, d'une \uncertain{paire} \uncertain{formée} des deux gén.\ de $G^{+}$ inverses l'un de l'autre, et d'un $G^{+}$-torseur $S$, est une équivalence de catégories.
\marginal{(au-dessus, renvoyé à la fin de la phrase) ou ce qui \uncertain{revient au même} \struck{$\mathbf{Z}$}-torseur}

\emph{Exer.} Exprimer $A$ \add{en termes des} données précédentes, ainsi que $R$ \uncertain{comme} \uncertain{partie} de $S \times R$\,…
\note{« $S \times R$ » : on attendrait plutôt une partie de $S \times A$ ; la lecture du dernier facteur est sûre sur la page.}

\page{62}
\marginal{trait vertical en marge gauche des deux premières lignes}
Cas des polygones réguliers, \struck{\ill{}} \uncertain{en} premier à la caractéristique ($k$ alg.\ clos)

$S$ ensemble des sommets, $A$ ens.\ des arêtes, $E$ espace affine ambiant.
Le barycentre $o$ de $S$ est fixé par le groupe $G$ des automorphismes de $\Pi$, de sorte que $E$ devient un vectoriel. Soit $G^{+}$ le sous-groupe des « rotations » de $\Pi$, $u$ un des 2 générateurs privilégiés, correspondant à une orientation \add{combinatoire} de $\Pi$, \struck{Considérons les deux valeurs propres de $u$, \ill{}} \add{soit $s_0 \in S$} [la donnée de $(s_0,u)$ équivaut à celle d'un repère combinatoire, i.e.\ d'un « épinglage combinatoire »] $\sigma$ l'automorphisme $\neq \mathrm{id}$ tel que $\sigma s_0 = s_0$, de sorte que
\[
\begin{cases}
\sigma^{2} = \mathrm{id}, \quad \sigma u \sigma^{-1} = u^{-1} \\
u^{n} = 1 \quad (n \text{ \uncertain{minimum}})
\end{cases}
\]
Soient $\zeta, \zeta'$ les valeurs propres de $u$. Si on avait $\zeta = \zeta'$, on aurait $u = \zeta(\mathrm{id}_E + v)$, $v$ \struck{nilpotent} \add{i.e.\ $v=0$}, et $u^{n} = \mathrm{id} \Rightarrow \zeta^{n} = 1$, $1 + nv = 0$ \uncertain{i.e.} $nv = 0$, d'où (comme $n \neq 0$ \uncertain{par \ill{}}) $v = 0$, donc $u = \zeta\,\mathrm{id}_E$, donc les $s_i = u^{i} s_0$ seraient sur une même droite, \uncertain{absurde}. Donc $\zeta \neq \zeta'$, soient $D_\zeta$, $D_{\zeta'}$ les droites propres. \struck{Mais} Comme $u \sim u^{-1}$, l'ensemble des

\page{63}
valeurs propres est stable par $z \mapsto z^{-1}$, i.e.\ \struck{$\zeta\zeta'$}
\[
\zeta\zeta' = 1 \quad\text{ou}\quad \zeta' = \zeta^{-1}.
\]
Plus précisément, comme $\sigma u \sigma^{-1} = u^{-1}$, $\sigma$ échange $D_\zeta$ et $D_{\zeta'}$, \struck{Choisissons $e_1 \in D_\zeta$ \ill{}, d'où} \add{si $s_0 = e_1 + e_2$ ($e_1 \in D_\zeta$, $e_2 \in D_{\zeta'}$) alors} $s_0 = \sigma s_0 = \sigma e_2 + \sigma e_1$ donc $\sigma e_1 = e_2$, $\sigma e_2 = e_1$ — on a aussi $e_1, e_2 \neq 0$ \uncertain{sinon $s_0 = 0$}, \uncertain{\ill{}} $e_1, e_2$ base de $E$ (définie \add{de façon unique}) \struck{\ill{}}.
\note{Le haut de la page est très raturé : une ligne entière biffée en zigzag (\struck{\ill{}}), une phrase intercalée sous la précédente (« \uncertain{indépendante} \ill{} $e_1 \mapsto e_2$ \ill{} \uncertain{épinglage} \ill{} ») et la fin de la parenthèse sont trop chargées pour être lues sûrement ; le sens retenu est que $(e_1, e_2)$ est la base de $E$ attachée à l'épinglage.}

Dans cette base, on a
\[
u = \begin{pmatrix} \zeta & 0 \\ 0 & \zeta^{-1} \end{pmatrix} \qquad
\sigma = \begin{pmatrix} 0 & 1 \\ 1 & 0 \end{pmatrix}
\]
\[
s_0 = \begin{pmatrix} 1 \\ 1 \end{pmatrix} \qquad
s_i = \begin{pmatrix} \zeta^{i} \\ \zeta^{-i} \end{pmatrix}
\]
On est dans le cas typique déjà étudié (sans avoir eu à supposer précédemment la \uncertain{car.}\ $0$, ni $\zeta$ d'ordre fini\,…). De plus, on voit que
\struck{\emph{Fonctions génératrices importantes} (ou \uncertain{avait} déjà remarqué \uncertain{précédemment} par \uncertain{rapport} à la base \uncertain{choisie})}
\note{Ces deux lignes sont barrées de traits obliques.}
\[
\alpha = \zeta + \zeta' = \zeta + \zeta^{-1} = \operatorname{Tr} u = \operatorname{Tr} u^{-1}
\]
est défini de façon intrinsèque en termes de $\Pi$, ou ce qui revient au même, la \uncertain{paire} \struck{\ill{}} \add{\uncertain{non ordonnée}}
\[
\{\zeta, \zeta'\} = \{\zeta, \zeta^{-1}\} = \{\zeta', \zeta'^{-1}\}
\]
d'éléments de $k^{*}$ inverses l'un de l'autre est définie intrinsèquement en termes de $\Pi$

\page{64}
On peut prendre pour $\zeta$ n'importe quelle racine de $1$, à l'exclusion de $\pm 1$ (racines d'ordre $1$ et $2$\,…).

Ainsi on trouve

\emph{Proposition} a) \struck{Tout \ill{}} polygone régulier épinglé à $n$ côtés, $3 \leq n < +\infty$ (donc $n$ fini), sur un corps alg.\ clos $k$ de car.\ $p$ premier à $n$, est isomorphe à un \add{des} polygones épinglés $\Pi_{\zeta,k}$ \add{dans $k^{2}$} \struck{$\Pi_{\zeta,k,\ill{},n}$} construits précédemment, \add{pour un $\Pi$,} \uncertain{où} $\zeta$ est défini modulo passage à $\zeta^{-1}$, \struck{\ill{}} \add{$\zeta$ est une racine primitive $n$-ième de $1$}.

b) On trouve ainsi une correspondance 1-1 entre classes d'isom.\ de polygones réguliers (ou \uncertain{les} classes d'isom.\ de polygones réguliers épinglés), entre \struck{\uncertain{couples}} \add{paires} $\{\zeta, \zeta^{-1}\}$ de racines primitives $n$-ièmes de $1$ inverses l'une de l'autre, et entre éléments $\alpha \in k$ tels que les $2$ solutions de l'équation
\[
\zeta^{2} - \alpha\zeta + 1 = 0
\]
soient des racines primitives $n$-ièmes de $1$.
\marginal{(en oblique, à gauche) \uncertain{les} $2$ \uncertain{racines} : \emph{invariants numériques} des polygones réguliers…}

c) Dans $\Pi_{\zeta,k}$, l'origine de $k^{2}$ est \uncertain{aussi} le barycentre de $S$ \add{$= \operatorname{Aut}(\Pi_{\zeta,k})$} \uncertain{et aussi} l'unique point fixe \struck{sous $G^{+}$} \add{sous $G$, ou $G^{+}$} \ill{} \add{l'unique \uncertain{point fixe}} \ill{} \uncertain{valeurs}
\note{La fin de la page est surchargée d'ajouts interlinéaires ; on ne lit sûrement que « l'unique point fixe sous $G$ », « $G^{+}$ », « $\{\zeta,\zeta^{-1}\}$ » et « $= \operatorname{Aut}(\Pi_{\zeta,k})$ ». L'énoncé enjambe sur la page suivante.}

\page{65}
propres de la partie homogène de \struck{\ill{}} \add{l'une} quelconque (ou \uncertain{des deux}) des deux générateurs privilégiés de $G^{+}$, les axes de coordonnées sont les droites propres \struck{\ill{}} \add{de} \uncertain{\ill{}} à l'épinglage près les valeurs propres \uncertain{et $\zeta^{-1}$ resp.}

\section{Formes invariantes (sous $\operatorname{Aut}(\Pi_{\zeta,n})$)}
\note{Titre souligné ; un mot biffé (\struck{\uncertain{quadratiques}}) précède « invariantes ».}

Si on regarde les fonctions polynômiales \struck{\ill{}} sur $E$ invariantes par $G$, comme l'origine est fixe, elles se décomposent en fonctions polynômiales homogènes, \struck{En degré $m$} \uncertain{ou formes}. L'invariance par $\sigma$ signifie la symétrie p.r.\ aux \struck{$x, y$} coordonnées \struck{$x$} $x$, $y$, donc
\[
\begin{cases}
F(x,y) = G(\sigma_1(x,y), \sigma_2(x,y)) \\
\sigma_1(x,y) = x + y, \quad \sigma_2(x,y) = xy
\end{cases}
\]
D'ailleurs, si
\[
F(x,y) = \sum_{\alpha+\beta = m} c_{\alpha\beta}\, x^{\alpha} y^{\beta}
\]
\note{un symbole biffé précède $F$.}
on a
\[
F(u(x,y)) = F(\zeta x, \zeta^{-1} y) = \sum c_{\alpha\beta}\, \zeta^{\alpha-\beta} x^{\alpha} y^{\beta}
\]
donc on a
\[
F \equiv F \circ u \iff c_{\alpha\beta} = 0 \text{ si } \beta - \alpha \not\equiv 0 \ (n)
\]

\page{66}
Donc les \struck{\ill{}} fonctions polynômiales invariantes ont une base sur $k$ formée de $1$ et des
\[
\begin{aligned}
\text{\struck{$\varphi_\alpha(x,y)$}}\qquad
\varphi_{\alpha,i}(x,y) &= x^{\alpha} y^{\alpha+in} + y^{\alpha} x^{\alpha+in} \\
&= x^{\alpha}y^{\alpha}(x^{in} + y^{in}) \\
&= (xy)^{\alpha}\bigl((x^{n})^{i} + (y^{n})^{i}\bigr)
\end{aligned}
\]
avec
\[
\alpha \in \mathbf{N}, \quad i \in \mathbf{N}^{*}
\]
\struck{Soit} Or si $X, Y$ sont des indéterminées, la théorie des fonctions symétriques nous dit que
\[
X^{i} + Y^{i} = S_i(\sigma_1(X,Y), \sigma_2(X,Y))
\]
\marginal{$\hookrightarrow$ polynôme universel : coeff.\ $\in \mathbf{Z}$ en $2$ variables $\sigma_1, \sigma_2$}
\[
= S_i(X+Y, XY)
\]
donc
\[
(x^{n})^{i} + (y^{n})^{i} = S_i(x^{n} + y^{n}, (xy)^{n})
\]
et par suite
\[
k[x,y]^{D} = k[x,y]^{\Delta_n} = k[\sigma_2, S_n]
\]
\note{Deux renvois au-dessus de la formule : l'exposant $D$ (suivi d'un indice biffé, \struck{$\zeta$}) est glosé « $= \operatorname{Aut}(\Pi_{\zeta,n})$ » ; $\Delta_n$ est glosé « groupe \uncertain{diédral}\ldots{} $\mathbf{Z}/n\mathbf{Z}$ (\ill{}) opérant sur $k^{2}$ de façon \uncertain{évidente} ».}
où
\[
\begin{cases}
\sigma_2(x,y) = xy \\
S_n(x,y) = x^{n} + y^{n}
\end{cases}
\]
\emph{NB} Les calculs marchent pour \uncertain{tout} $n \in \mathbf{N}^{*}$ (pas besoin que $n \geq 3$) et sur \uncertain{tout} anneau de \uncertain{base} — le meilleur des \uncertain{trucs} étant $\mathbf{Z}$\,! Si $n$ est infini, il faut revenir au cas \uncertain{d'un} $\zeta$

\page{67}
donné, \uncertain{solution} d'une \uncertain{équation} \uncertain{structurelle} d'ordre infini sur \uncertain{chaque} \uncertain{fibre}, on trouve \uncertain{\ill{}} comme anneau d'invariants $k[\sigma_2]$\,…

\emph{Corollaire} Les fonctions polynômiales de degré $\leq 2$ invariantes par $D_\zeta = \operatorname{Aut}(\Pi_{\zeta,n})$ sont les fonctions
\[
axy + b \qquad (a, b \in k)
\]
Donc à scalaire près une unique forme \add{$q(x,y) =$ \struck{\ill{}} $\sigma_2(x,y) = xy$} quadratique invariante, à savoir par normalisation $\sigma_2(s) = 1$ si $s \in S$. Une unique \uncertain{conique} « affine » circonscrite, i.e.\ \uncertain{contenant} les sommets, à savoir celle d'équation
\[
xy - 1 = 0
\]
\struck{\ill{}}
\marginal{(en oblique, à gauche) \uncertain{Kif}\,! \uncertain{Ne marche pas que si} $n \geq 5$, avec \uncertain{\ill{}} \uncertain{l'interprétation} \ill{} conique \uncertain{\ill{}}. Ici on a \uncertain{\ill{}} une conique circonscrite canonique — d'équation invariante par $G$}
\struck{\ill{}}

Pour un polygone régulier quelconque (sur un corps $k$ alg.\ clos de car.\ $p$ premier à $n$) \add{\uncertain{à scalaire près}} \uncertain{non} épinglé — il y a donc une unique forme quadratique invariante $q_0$, \uncertain{normalisée} par $q_0(s) = 1$ ($s \in S$)

\page{68}
et une unique \struck{\uncertain{quadrique}} \add{conique} circonscrite à $\Pi$ i.e.\ à $S$, savoir la conique
\[
C \quad\text{d'équation}\quad q_0(x) - 1 = 0
\]
Notons que $q_0$, considérée comme forme quadratique sur l'espace des translations de $E$, fait de $E$ un « plan euclidien » — la forme est en effet non dégénérée. Le carré de la longueur des côtés du polygone $\Pi$ est donné par
\[
\begin{aligned}
q_0(s_1 - s_0) &= q_0(\zeta - 1, \zeta^{-1} - 1) = (\zeta-1)(\zeta^{-1}-1) \\
&= 2 - \alpha
\end{aligned}
\]
i.e.
\[
q_0(a) = 2 - \alpha
\]
\note{ici $a$ désigne un côté (arête) de $\Pi$.}
On notera que c'est \uncertain{toujours} $\neq 0$ i.e.\ \struck{\ill{}} $\alpha \neq 2$, puisque $\zeta \neq 1$. [On pourrait donc normaliser \struck{\ill{}} $q$ \uncertain{différemment}, en normalisant par $q_0(a) = 1$. Dans le cas général, c'est cette normalisation-là qui sera la meilleure, comme nous verrons\,…]

\struck{\emph{Commentaires « philosophiques »} On a trouvé, pour le groupe $D_n$ fixé, et \uncertain{des} corps}
\note{Ces deux lignes sont barrées de traits obliques, ainsi que le haut de la page suivante jusqu'à « \ill{} près. ».}

\page{69}
\struck{\uncertain{\ill{}} des variétés $k$ \add{alg.\ clos (pour simplifier)} de car.\ première à $n$, \struck{\ill{}} munis d'un $\alpha \in k$ tel que les solutions de $\zeta^{2} - \alpha\zeta + 1 = 0$ sont des racines primitives $n$-ièmes de $1$ (ce qui peut s'exprimer par la condition $F_n(\zeta) = 0$\,…) une famille \add{de $D_n$} \struck{\ill{}} de représentations affines de dim $2$ \add{des isomorphismes} \struck{représentations polygones réguliers} — ou encore celle \uncertain{associée} \ldots{} \uncertain{\ill{}} des couples formés d'un polygone régulier d'isomorphisme, et d'un épinglage de celui-ci — ou encore des \uncertain{homomorphismes} $D_n \to \operatorname{Aff}(\mathbf{3}k)$ définis à automorphisme intérieur près.}
\note{« $\operatorname{Aff}(\mathbf{3}k)$ » : lecture incertaine du dernier symbole, peut-être $\operatorname{Aff}(E_k)$ ou $\operatorname{Aff}_2 k$.}

\emph{Remarque} Si on s'était borné à exiger, pour un polyèdre \add{affine} \struck{affine} \uncertain{régulier}, que le groupe de ses automorphismes \emph{affines} soit transitif sur l'ens.\ des sommets \add{(sans exiger rien \struck{\ill{}} sur les autres \uncertain{ensembles})}, on aurait trouvé beaucoup plus de choses. Par exemple, prenant \add{le groupe $G^{+}$} \add{\uncertain{des} rotations} combinatoires, provenant d'un groupe de transformations affines, on aurait eu à choisir des $u$ tels que $u^{n} = 1$ (sans \uncertain{se} \uncertain{donner} de $\sigma$ tel que $\sigma u \sigma^{-1} = u^{-1}$) — ce qui revient \uncertain{à se} donner deux v.p.\ $\zeta, \zeta'$ racines \struck{\uncertain{pr}} $n$-ièmes de $1$, d'ordre $d, d'$, avec

\page{70}
$n = \operatorname{ppcm}(d, d')$ (qui exprime que $u$ est d'ordre $n$ exactement) — $S$ se définissant comme l'orbite d'un $s_0 \in E \setminus (D_\zeta \cup D_{\zeta'})$, \struck{et $A$} \ldots{} l'ens.\ des $s_i = u^{i} s_0$ \struck{et} ($i \in \mathbf{Z}/n\mathbf{Z}$), et $A$ est formé des droites \add{$a_i = \operatorname{dr}\{s_i, s_{i+1}\}$} (il faut encore vérifier, en exigeant que $a_i$ ne soit incident qu'à $s_0$ et $s_1$, cette condition est laissée en exercice\,…)

Par contre, si on exige \add{de plus} qu'il existe une \uncertain{métrique} \add{$\delta$} \uncertain{non} dégénérée invariante par les \struck{$\delta$} $\in G^{+}$, et en insistant pour que les automorphismes soient \uncertain{notés} des \emph{rotations} (comme me l'avait suggéré Olivier Leborgne), on sauve la mise. En effet, les rotations sont unimodulaires, donc on trouve pour les valeurs propres de $u$ \uncertain{une} $\{\zeta, \zeta^{-1}\}$, et on retrouve les cas précédents\,… Je ne sais pas \uncertain{si} \uncertain{on simplifie} \uncertain{\ill{}} encore en \emph{exigeant} simplement que les $\in G^{+}$ respectent la métrique, sans plus\,…
\note{« Olivier Leborgne » : lecture incertaine du nom propre.}

\page{71}
Fonctions polynômiales rel.\ invariantes sous $G^{+}$, sous $G$.

\struck{\ill{}} Les $x^{i}y^{j}$ forment une base de $\operatorname{Sym}^{*}(\check{E})$, formée de vecteurs propres pour $u$, de valeurs propres $\zeta^{i-j}$. \struck{\ill{}} Les \uncertain{val.}\ des valeurs propres de $G^{+} \simeq \mathbf{Z}/n\mathbf{Z}$ dans les valeurs propres de \uncertain{$G^{+}$} opérant sur $\operatorname{Sym}^{*}(\check{E}) \simeq$ Fonctions polyn.\ $(E)$ sont donc les $\zeta^{d}$ ($d \in \mathbf{Z}/n\mathbf{Z}$), \struck{et les} fonctions correspondant à cette valeur propre sont les combinaisons linéaires des \struck{$x^{j+d}y^{j} = (xy)^{j}x^{d}$} \ldots{} i.e.
\[
x^{i}y^{j} \qquad i - j \equiv d \ (n)
\]
i.e.\ les
\[
x^{j+d+nk}\,y^{j} =
\begin{cases}
(xy)^{j}\, x^{d+nk} & \text{si } d + nk \geq 0 \\
(xy)^{\uncertain{j'}}\, y^{-d+nk'} & \text{si } \uncertain{-d+nk' \geq 0}
\end{cases}
\]
\note{La fin des deux conditions mord sur le bord de la feuille ; la seconde ligne de l'accolade est en partie surchargée.}

\struck{\uncertain{Cas $d=0$}} En degré $0$ on obtient \uncertain{donc} \uncertain{le} \ldots{}
\[
\operatorname{Sym}^{*}(\check{E})^{G^{+}} \simeq k[xy, x^{n}, y^{n}]
\qquad (U_n V_n - W_n^{n} = 0)
\]
\note{Au-dessus des générateurs, renvois $W = W_n = xy$, $U = U_n = x^{n}$, $V = V_n = y^{n}$. La relation entre parenthèses mord sur le bord ; on lit $U_nV_n - W_n^{n} = $ \ill{}.}
\struck{(\ill{}) par $x^{d}$ et $y^{-d+nk}$, \ill{}}
c'est engendré, \uncertain{si} $0 < d < n$, par
\[
x^{d} = X_d \quad\text{et}\quad y^{n-d} = Y_d,
\]
avec les conditions
\[
V X_d - (W)^{d} Y_d = 0 \quad \ldots
\]
Pour les fonctions polynômiales de degré $\leq 2$,

\page{72}
ayant comme base
\[
1,\ x,\ y,\ x^{2},\ xy,\ y^{2}
\]
les valeurs propres de $u$ correspondantes sont resp.
\[
1,\ \zeta,\ \zeta^{-1},\ \zeta^{2},\ 1,\ \zeta^{-2}
\]
donc les valeurs propres qui y figurent sont \struck{\uncertain{respectivement}} parmi les \add{$5$} valeurs
\[
\zeta^{-2},\ \zeta^{-1},\ \zeta^{0} = 1,\ \zeta^{1} = \zeta,\ \zeta^{2}.
\]
Celles-ci sont \uncertain{toutes} distinctes \uncertain{si} $n \geq 5$, dans ce cas les \struck{fonctions} \add{p.\ \uncertain{homogènes}} \uncertain{semi-invariantes} \add{sous $G^{+}$} de degré $\leq 2$ \struck{\ill{}} sont exactement les fonctions
\[
\begin{aligned}
&a + bxy && \text{valeur propre } \zeta^{0} = 1 \\
&ax \ (\text{\uncertain{resp.}}\ ay) && \text{valeur propre } \zeta\ (\text{resp.\ } \zeta^{-1}) \\
&ax^{2}\ (\text{resp.\ } ay^{2}) && \text{valeur propre } \zeta^{2}\ (\text{resp.\ } \zeta^{-2})
\end{aligned}
\]
Si on veut de plus la quasi-invariance sous $\sigma$ \add{i.e.\ sous $G$}, \struck{\ill{}} $G \simeq D_n$ tout entier, notant que $\sigma(ax) = ay$, $\sigma(ay) = ax$, $\sigma(ax^{2}) = ay^{2}$, $\sigma(ay^{2}) = ax^{2}$, on trouve que les seules fonctions \add{p.\ \uncertain{homogènes} de degré $\leq 2$} quasi-invariantes sous $G$ sont les $a + bxy$ (qui sont invariantes)

Donc les seules coniques invariantes sous $G$ sont les coniques d'équation $a + bxy$, donc si on veut les coniques circonscrites à $\Pi$ il faut $a + b = 0$ i.e.\ $b = -a$, on trouve la \struck{\uncertain{seule}} conique \emph{invariante} circonscrite, d'équation $xy - 1 = 0$\,…
\note{« $a + b = 0$ » : la condition écrite est sans doute à lire $a + b = 0$ avec la normalisation $xy(s) = 1$ des sommets.}

\page{73}
En fait, le choix des « coordonnées épinglées universelles » montre que c'est la seule conique circonscrite (celles en \uncertain{dimension} \uncertain{autres} \uncertain{triangles} invariantes\,…)

\emph{Cas $n = 3$} i.e.\ $\zeta^{2} = \zeta^{-1}$, $\zeta^{-2} = \zeta$, donc les fonctions p.l.\ quasi-inv.\ \add{sous $G^{+}$ \uncertain{i.e.\ sous $u$}} de degré $\leq 2$ sont les
\[
\begin{cases}
a + bxy & \text{v.p.\ } 1 \\
ax + by^{2} & \text{v.p.\ } \zeta \\
ay + bx^{2} & \text{v.p.\ } \zeta^{-1}
\end{cases}
\]
\note{les deux dernières lignes portent un trait de répétition pour « v.p. » ; dans la troisième, $ay^{2}$ est corrigé en $ay$.}
mais parmi celles-ci, les quasi-inv.\ sous $G$ sont encore les $a + bxy$.

\emph{Cas $n = 4$} i.e.\ $\zeta^{2} = \zeta^{-2}$, les fonctions p.l.\ de degré $\leq 2$ quasi-inv.\ sous $G^{+}$ \add{i.e.\ sous $u$} sont
\[
\begin{cases}
a + bxy & \text{v.p.\ } 1 \\
ax \ (\text{\uncertain{resp.}}\ ay) & \text{v.p.\ } \zeta\ (\zeta^{-1}) \\
ax^{2} + by^{2} & \text{v.p.\ } \zeta^{2} = \zeta^{-2}
\end{cases}
\]
Parmi celles-ci, les fonctions quasi-inv.\ sous $G$ sont celles qui sont les $a + bxy$ et \uncertain{distinctes} \struck{telles qu'il y a deux} \add{qui forment} \uncertain{la} famille : $ax^{2} + by^{2}$, \ldots{} une \uncertain{par} famille de coniques inv.\ sous $G$. Dans chacune, il y a \struck{exactement} une conique circonscrite à $S$, savoir \add{l'habituelle $xy - 1 = 0$}, \add{et \uncertain{de plus} \ldots{}} $x^{2} - y^{2} = 0$

\page{74}
\section{Réalisation géométrique des polyèdres \struck{\uncertain{affines}} combinatoires en coordonnées affines universelles.}

1.\ Programme de travail.

\struck{\ill{}} 1.\ Soit $\Pi$ polyèdre combinatoire donné; on se propose d'étudier ses réalisations géométriques \add{\uncertain{régulières}} dans des \struck{\uncertain{espaces}} \add{plans} affines variables $E$. Notons que si deux telles réalisations sont $\Pi$-isom., il y a un $\Pi$-isom.\ \emph{unique} entre elles, par définition. Donc dans \uncertain{les} classes d'isomorphie (\uncertain{fixant} les idées), il y a un \uncertain{certain} \uncertain{nombre} de \add{$(E_\alpha, \varphi_{\alpha s}, \varphi_{\alpha a})$} représentants, définis à isom.\ unique près. \uncertain{En} particulier, le plan affine \add{$E_\alpha$} de cette réalisation \uncertain{géométrique} est défini à isom.\ \emph{unique} près.

On indexera $\alpha$ à l'aide d'invariants numériques (\uncertain{essentiellement} $\alpha$ — ensemble \add{« \uncertain{\ill{}} numérique »} de \uncertain{scalaires}, \uncertain{comme} $\zeta + \zeta^{-1}$ défini universellement précédemment) et pour bien faire il faudra \uncertain{fournir} \add{(à isom.\ unique près)} explicites $E_\alpha$, $\varphi_{\alpha s}$, $\varphi_{\alpha a}$ à l'aide de $\Pi$ et de cet \uncertain{invariant} numérique.

On ne peut \uncertain{commencer} se donner un repère \add{$r \in \operatorname{Rep}(\Pi)$} pour épingler $\Pi$, de sorte que $\Pi$ devient

\page{75}
canonit.\ isom.\ à un polygone combinatoire-type $\Pi_n$ ($3 \leq n \leq +\infty$). D'ailleurs tous ces $\Pi_n$ sont \add{canoniquement} quotients du même $\Pi_\infty$ (si $n | m$, $\Pi_n$ est quotient de $\Pi_m$), par cet \add{\uncertain{l'image} \uncertain{du} rep.\ de $\Pi_n$} l'image du repère $r_\infty$ de $\Pi_\infty$, donc la cat.\ (discrète) des repr.\ affines régulières de $\Pi_n$ en plans (comme sous-catégorie pleine) de celle des repr.\ régulières de $\Pi_\infty$. Dans les \uncertain{plans} pratiques nous nous en ferons l'étude en trois étapes :

a) Étudier les réal.\ géom.\ rég.\ de $\Pi_\infty$ sur un corps (plus gén.\ \uncertain{un anneau} \ldots{} \uncertain{(\ill{})}) \struck{\ill{}} quelconque. On verra entre autres qu'elles sont caractérisées, à isom.\ près, par un invariant numérique $\alpha \in k$, qui peut être choisi arbitrairement.

b) Expliciter les conditions sur $\alpha$ pour que la réal.\ géom.\ corresp.\ \uncertain{qui} de $\Pi_\infty$ se factorise par $\Pi_n$, et plus

\page{76}
précisément pour que la \struck{\uncertain{factorisation}} \add{réalisation} géom.\ obtenue de $\Pi_n$ ($3 \leq n \leq \infty$) soit fidèle. De cette façon, on trouvera une partition ou classification des \struck{\uncertain{formes}} $n$-polygones réguliers sur $k$, indépendamment de la donnée d'un $\Pi$-épinglage, par \struck{\ill{}} un invariant $\alpha \in k$ satisfaisant certaines conditions.

c) Déduire de ceci les \struck{\ill{}} propriétés essentielles \struck{des} (en particulier la classification) des $n$-polygones réguliers, quel que soit $n$, et indépendamment de tout épinglage par un repère. (Cette étape est essentiellement triviale\,…)

\section{2. Reformulation du pb des réalisations géom.\ régulières de $\Pi_\infty$ dans des plans affines $E$ (sur un corps $k$, pour simplifier\,…)}
\note{Titre marqué d'un trait vertical en marge gauche.}

Rotations ($G = \operatorname{Aut}(\Pi_\infty) \simeq D_\infty$ \add{$= G_{0,1}$} $= \{\sigma_0, \sigma_1 \mid \sigma_0^{2} = \sigma_1^{2} = 1\}$)

a) Représentation
\[
(1) \qquad \varphi_G : D_\infty \longrightarrow \operatorname{Aff}(E)
\]

\page{77}
i.e.
\[
(2) \qquad \boxed{\sigma_0, \sigma_1 \in \operatorname{Aff}(E) \qquad \sigma_0^{2} = \sigma_1^{2} = \mathrm{id}_E}
\]
(b)
\[
(3) \qquad \boxed{r = (s_0, a_0) \in E \times \operatorname{Dr}(E), \quad (s_0, a_0) \text{ incidents}}
\]
\emph{NB} Ces deux données définissent $\varphi_S$, $\varphi_A$ par
\[
(4) \qquad
\begin{cases}
\varphi_S(\underline{s}_n) = s_n = u^{n} s_0 \\
\varphi_A(\underline{a}_n) = a_n = u^{n} a_0
\end{cases}
\]
avec
\[
(5) \qquad u = \sigma_0 \sigma_1
\]
\struck{Mais on} On trouve de \uncertain{même}, pour $\underline{r} \in \operatorname{Rep}(\Pi_\infty)$
\[
\begin{aligned}
&\underline{r} = \underline{g}\,\underline{r}_\infty \qquad (\underline{g} \in \operatorname{Aut}(\Pi_\infty) = D_\infty) \\
(5) \qquad &\varphi_{\mathrm{rep}}(\underline{r}) = \varphi_{\mathrm{rep}}(\underline{g}\,\underline{r}_\infty) = g\,\varphi_{\mathrm{rep}}(\underline{r}_\infty) = g.r
\end{aligned}
\]
\note{Deux formules portent le numéro (5) sur la page ; le second chiffre est peut-être un (6) mal formé, la suivante étant aussi numérotée (6).}
i.e.
\[
(6) \qquad \varphi(\underline{g}\,\underline{r}_\infty) = g.r \quad\text{si}\quad g = \varphi_G\,\underline{g}
\]
et les formules (4), \struck{\ill{}} peuvent être considérées comme déduites de (6) par passage au quotient via
\note{diagramme redessiné : $\operatorname{Rep}(\Pi_\infty)$ en haut, deux flèches descendantes vers $S(\Pi_\infty)$ et $A(\Pi_\infty)$, étiquetées $\underline{\sigma}_1$ et $\underline{\sigma}_0$ comme sur la page.}

\begin{tikzcd}
& \operatorname{Rep}(\Pi_\infty) \arrow[dl, "\underline{\sigma}_1"'] \arrow[dr, "\underline{\sigma}_0"] & \\
S(\Pi_\infty) & & A(\Pi_\infty)
\end{tikzcd}

i.e.\ on doit avoir $\varphi_S$, $\varphi_A$ définis par

\page{78}
commutatifs dans

\begin{tikzcd}
\operatorname{Rep}(\Pi_\infty) \arrow[r, "\varphi_{\mathrm{rep}}"] \arrow[d, "\mathrm{pr}_1"] \arrow[dd, bend right=50, "\mathrm{pr}_2"'] & \operatorname{Drap.aff}(E) \arrow[d, "\mathrm{pr}_1"] \arrow[dd, bend left=50, "\mathrm{pr}_2"] \\
S(\Pi_\infty) \arrow[r, "\varphi_S"] & E \\
A(\Pi_\infty) \arrow[r, "\varphi_A"] & \operatorname{Dr}(E)
\end{tikzcd}

\note{Diagramme (7) redessiné ; sur la page la flèche $\mathrm{pr}_2$ de droite part de $\operatorname{Drap.aff}(E)$ en croisant la flèche $\varphi_S$. « $\operatorname{Drap.aff}(E)$ » : lecture probable, les drapeaux affines (point, droite incidente) de $E$.}

La condition [pour $\varphi_G$ \add{(cf (1), (2))} et le $G$-hom.\ $\varphi_{\mathrm{rep}} : \operatorname{Rep}(\Pi_\infty) \to \operatorname{Drap.aff}(E)$ donnés, cette dernière donnée correspondant à la donnée \add{(des $\varphi_S$ et $\varphi_A$)} de (4)] pour que $\mathrm{pr}_i \circ \varphi_{\mathrm{rep}}$ passe au quotient — ou encore pour que $\varphi_S$ et $\varphi_A$ soient compatibles avec les relations d'incidence, c'est que l'on ait
\[
\sigma_1\, \mathrm{pr}_1(r) = \mathrm{pr}_1(r), \quad \sigma_0\, \mathrm{pr}_2(r) = \mathrm{pr}_2(r)
\]
i.e.
\[
(8) \qquad \boxed{\sigma_1 s_0 = s_0, \quad \sigma_0 a_0 = a_0}
\]
Il reste à exprimer que les $\uncertain{\ill{}}$ engendrent affinement $E$. On va introduire
\note{la suite de l'argument passe à la page suivante.}

\page{79}
\[
s_0, \quad s_1 = \sigma_0 s_0\ \bigl(= u s_0 = \sigma_0 \underbrace{\sigma_1 s_0}_{s_0}\bigr), \quad
s_{-1} = \sigma_1 s_1\ (= s_{-1} = \sigma_1 \sigma_0 s_0)
\]
\note{le dernier terme est écrit d'abord $\underline{s}_1'$, surchargé en $s_{-1}$ ; la parenthèse finale mord sur le bord de la feuille.}
On a la

\emph{Lemme} \uncertain{Étant} donnés (2) (3) satisfaisant (8), les conditions suivantes sont équivalentes :
\begin{quote}
(a) Les données correspondent à une réalisation géom.\ de $\Pi_\infty$, i.e.\ les $s_n = u^{n} s_0 = (\sigma_0\sigma_1)^{n} s_0$ sont affinement indépendants

(b) $s_0$, $s_1$ et $s_{-1}$ sont affinement indépendants.

(c) Les applications $\varphi_S$ et $\varphi_A$ de (3) sont l'une et l'autre non constantes.
\end{quote}
\note{Les trois conditions sont réunies par un trait vertical en marge.}
\note{En (a), « affinement indépendants » pour une suite infinie de points d'un plan ne peut s'entendre au sens littéral ; on transcrit tel quel.}

En effet, b) $\Rightarrow$ (a) $\Rightarrow$ (c) sont triviaux, prouvons c) $\Rightarrow$ a). On suppose c), il faut prouver

1°) $s_1 \neq s_0$ sinon $s_0$ serait fixé par $\sigma_1$ et $\sigma_0$, donc par $G$, donc par les $u^{n}$, donc $\varphi_S$ serait constante.

2°) $s_{-1} \notin \operatorname{dr}(s_0, s_1)$. En effet, \struck{\ill{}}
\struck{Comme, $\sigma_0$ \ill{}}
\[
(9) \qquad a_0 = \operatorname{dr}(s_0, s_1)
\]
Car $s_0 \in a_0$, et comme $\sigma_0 a_0 = a_0$ \add{\uncertain{par (8)}}, on a aussi $\sigma_0 s_0 = s_1 \in a_0$. \struck{Comme} \uncertain{Or} \ldots
\[
\sigma_0 a_0 = \operatorname{dr}(\sigma_0 s_0 = s_1,\ \sigma_1 s_1 = s_{-1})
\]
et si on avait $s_{-1} \in a_1$, on aurait $\sigma_1 a_0 = a_0$
\note{L'indice de ce dernier $\sigma$ mord sur le bord de la feuille et se lit mal ; $\sigma_1$ est demandé par la suite. La formule $\sigma_0 a_0 = \operatorname{dr}(\ldots)$ est écrite telle quelle ; on attendrait $\sigma_1 a_0$. L'argument continue au-delà du lot.}

\page{80}
donc $a_0$ serait fixée par $\sigma_0$ et $\sigma_1$, donc par $G$, donc par les $u^{n}$, donc $\varphi_A$ serait constante, cqfd.
En même temps, on a retrouvé (pour mémoire) (9), qui montre que les données (2) (3) satisfaisant (8) + condition \uncertain{(b)} \uncertain{équivalent} à : (2), \struck{\ill{}} \add{$+$ donnée} (\uncertain{seulement} $s_0$) satisfaisant que $s_0, s_1$ et $s_{-1}$ \ldots{}

\emph{Proposition} Les réalisations géom.\ régulières de $\Pi_\infty$ dans $E$ correspondent 1-1 aux triples $(\sigma_0, \sigma_1, s_0)$, où
\[
(10) \qquad
\begin{cases}
\sigma_0, \sigma_1 \in \operatorname{Aff}(E) \\
s_0 \in E
\end{cases}
\]
satisfaisant les conditions

a) $\sigma_0^{2} = \sigma_1^{2} = \mathrm{id}$

b) $\sigma_1 s_0 = s_0$ i.e.\ $s_0 \in E^{\sigma_1}$

c) $s_0$, $\sigma_0 s_0 = s_1$, et $s_{-1} \overset{\mathrm{def}}{=} \sigma_1 s_1$ forment une base affine de $E$.

La repr.\ $\varphi_G : G = \operatorname{Aut}(\Pi_\infty) \simeq G_{0,1}$ associée est définie par
\[
(11) \qquad \varphi_G(\underline{\sigma}_0) = \sigma_0, \quad \varphi_G(\underline{\sigma}_1) = \sigma_1
\]
et $\varphi_S$ et $\varphi_A$ sont données par
\[
\left.
\begin{aligned}
(12) \qquad & \varphi_S(\underline{s}_n) \overset{\mathrm{def}}{=} s_n = u^{n} s_0 \\
(13) \qquad & \varphi_A(\underline{a}_n) = \operatorname{dr}(s_n, s_{n+1}) = u^{n} \operatorname{dr}(s_0, s_1)
\end{aligned}
\right| \; n \in \mathbf{Z}
\]
où
\[
(14) \qquad u = \sigma_0 \sigma_1
\]
\note{La page s'arrête ici ; l'exposé (« 2. Reformulation… ») se poursuit au-delà du lot.}

\end{document}
